Large language models trained or fine-tuned on private data can memorise parts
of that data~\citep{carlini2021extracting,shokri2017membership}.
Membership inference attacks (MIA) try to detect this: given a text and some
access to a model, can an attacker tell whether that text was in the training set?
This is the standard empirical test for privacy leakage, and it matters practically
whenever models are fine-tuned on sensitive corpora like medical records, legal
documents, or private codebases~\citep{carlini2022membership}.

Most MIA work targets autoregressive (AR) models, where a per-token log-probability
is directly available and serves as a natural score~\citep{duan2024mimir,yeom2018privacy}.
Masked Diffusion Language Models~\citep{austin2021d3pm,sahoo2024mdlm,nieLLaDA2025}
work differently: they learn to reconstruct randomly masked tokens rather than
predict the next one. There is no single unconditional likelihood to read off.
Instead, the loss depends on which tokens are masked and at what ratio, and this
dependency turns out to carry membership information that scalar-based attacks miss.

We focus on the ELBO trajectory: how the masked reconstruction loss changes as
the masking ratio $\alpha$ increases from 0.05 to 0.50.
A model that has memorised a text reconstructs it cheaply at any masking level,
producing a lower and more concave loss curve than for unseen text.
Collecting these values at four masking levels and adding a few supporting signals
(predictive entropy, hidden-state norms) gives a 46-dimensional feature vector
that a gradient-boosted classifier can use to separate members from non-members.

Prior work on MDLMs includes SAMA~\citep{shi2025sama}, a grey-box attack that
aggregates NLL scores across random token subsets. Our white-box trajectory
features beat it by +0.062 AUC on average. For autoregressive models,
\citet{duan2024mimir} show that standard loss-based attacks largely fail on
pretrained models but remain viable after fine-tuning.
MIA for diffusion image models is studied in~\citet{dockhorn2022dpdm}.
No prior work applies white-box trajectory extraction to discrete diffusion LMs.

We make four contributions:

(1) \textbf{White-box trajectory MIA.} We extract a 46-dim feature vector at four masking
ratios and train XGBoost, LightGBM, and MLP classifiers via 5-fold CV.
Mean AUC of 0.878 across six MIMIR domains, peaking at 0.930.

(2) \textbf{Shadow-model transfer.} We train classifiers on $K=3$ shadow MDLMs
fine-tuned on surrogate data from different domains, then apply them directly
to the target domain with no retraining. This achieves 0.858 mean AUC with no
target-domain labels, closing within 0.020 of the white-box oracle.
Shadow transfer is a practical, near-equally effective attack path.

(3) \textbf{Feature ablation.} A LOSO analysis across all six domains isolates
the ELBO trajectory as the primary signal (mean drop 0.142 when removed).
All four attention-based signal groups contribute less than 0.003 each.

(4) \textbf{Attention features do not transfer.} When only attention features
are used in the shadow-model setting, AUC drops to 0.525, confirming they are
domain-specific and not reusable across domains. ELBO-based signals are.
